% What we need from Chapter 3 --- spec \S4.1 row 2.
% Boxed formulas from Project_specification.md \S6; pointers to Chapter 3;
% no proofs beyond one-liners. All correction-history remarks deleted.

\section{What we need from Chapter 3}
\label{sec:recap}

Chapter~3 defined the birth--death--catastrophe process and derived its
principal single-cell quantities in closed form. This chapter builds on that
foundation, and everything it needs is restated here, boxed, so that what
follows is self-contained. Derivations are at most sketched; the full
accounts live in Chapter~3, to which each box points.
Section~\ref{sec:notation-tab} fixes the notation used throughout; the
remaining subsections collect the closed forms that will be quoted
repeatedly.

\subsection{Canonical notation}
\label{sec:notation-tab}

The notation of this chapter is collected in Table~\ref{tab:notation}; it is
the notation of Chapter~3, unchanged. When, in Chapter~4b, the single-cell
theory is joined to a population model, further symbols are introduced
there, chosen to avoid collisions with these.

\begin{table}[H]
\centering
\renewcommand{\arraystretch}{1.25}
\small
\begin{tabular}{@{}llp{0.44\textwidth}@{}}
\toprule
Symbol & Meaning & Notes \\
\midrule
$\lambda,\mu,\delta$ & birth, death, catastrophe rates (per capita) & \\
$\xt$ & intracellular count (CTMC) & \\
$H$ & absorbing rupture state & \\
$\wt$ & released count from one cell & $0$ until burst \\
$\tau$ & burst time (possibly $\infty$) & \\
$\Kb$ & burst size (random variable) & never bare $K$, since $K(t)=\mean{\xt^2}$ \\
$I(t)$ & $\Prob{\text{no burst by }t}$ & \\
$D(t)$ & $\Prob{\text{cleared internally, no burst}}$ & $= p_0(t)$ \\
$\Ifix(t)$ & $\Prob{\xt\notin\{0,H\}} = I(t)-D(t)$ & productive infection \\
$J(t),K(t)$ & $\mean{\xt}$, $\mean{\xt^2}$ & \\
$V(t)$ & $\mean{\wt}$, cumulative mean release & \\
$a,b$ & characteristic roots, $b<1<a$ & \\
$A,B,\theta$ & $a-1$, $1-b$, $\lambda(a-b)$ & \\
$\kappa$ & $1+\delta/(2\lambda)$ & appears in $K(t)$ and $K(I)$ \\
$w$ & $\ee^{\theta t}$ & the recurring exponential \\
\bottomrule
\end{tabular}
\caption{Canonical notation for this chapter (Chapter~3's regime).}
\label{tab:notation}
\end{table}

\subsection{Process definition and the joint process}

One cell, one intracellular population. The count $\xt$ is a continuous-time
Markov chain on $\{0,1,2,\ldots\}\cup\{H\}$ which, from any state $n\ge1$,
jumps
\begin{align*}
n &\to n+1 &&\text{at rate } \lambda n \quad\text{(birth)},\\
n &\to n-1 &&\text{at rate } \mu n \quad\text{(death)},\\
n &\to H &&\text{at rate } \delta n \quad\text{(catastrophe)}.
\end{align*}
Both $0$ and $H$ are absorbing: the population either clears itself
internally, or the cell ruptures. Unless stated otherwise we start from a
single founder, $X_0=1$. The process of classical references
\cite{brockwell1982birth,karlin1982linear,di2008note} is slightly more
general (immigration, distributed catastrophe sizes); the specialisation
above is the one with the closed-form family we need.

\figureflag{F4a.1}{Joint-process schematic}{%
\textbf{Purpose:} display the object of this chapter --- the joint process
$(X_t,W_t)$ --- at a glance: the stochastic growth of the load, the burst at
$\tau$, and the jump of the released count from $0$ to $X_{\tau^-}$.
\textbf{Type:} TikZ schematic, inline with chapter fonts. \textbf{Content:}
a horizontal time axis; one representative sample path of $X_t$ drawn as a
step function starting at $X_0=1$, wandering up and down through states
$1,2,\ldots$, terminating in a vertical jump to a marked state $H$ at time
$\tau$; below it, on a second axis, the corresponding path of $W_t$: zero
until $\tau$, then a vertical jump to the level $X_{\tau^-}$ and a flat
continuation thereafter; annotate $\tau$, $X_{\tau^-}$, and label the state
$H$ with an ellipse as in Chapter~3's transition diagram. \textbf{Labels:}
$t$ on the horizontal axis; $X_t$ and $W_t$ on the two vertical axes;
annotations $\tau$, $X_{\tau^-}$, $H$. \textbf{Style:} TikZ, inline with
chapter fonts, black paths, the burst jump optionally emphasised in colour.
\textbf{Placement:} \Secref{sec:recap}, in the ``Process definition and the
joint process'' subsection, after the transition-rate display.}

The rupture event is what matters for applications, so we record the release
count alongside the load. Let $\tau$ be the time of catastrophe (possibly
$\tau=\infty$ if the population clears internally first). The released count
$\wt$ is
\[
\wt = \begin{cases}
0, & t<\tau,\\
X_{\tau^-}, & t\ge\tau,
\end{cases}
\]
that is, at rupture the whole pre-burst load is dumped and then held. All
single-cell release quantities in this chapter are functionals of the joint
process $(\xt,\wt)$.

\subsection{Roots and the workhorse identity}

Write
\begin{equation}
\eta=\frac{\lambda+\mu+\delta}{2\lambda},\qquad
a=\eta+\sqrt{\eta^2-\tfrac{\mu}{\lambda}},\qquad
b=\eta-\sqrt{\eta^2-\tfrac{\mu}{\lambda}},
\end{equation}
the two roots of the quadratic $\lambda f^2-(\lambda+\mu+\delta)f+\mu$, and
set
\begin{equation}
A:=a-1>0,\qquad B:=1-b>0,\qquad \theta:=\lambda(a-b)
=\sqrt{(\lambda+\mu+\delta)^2-4\lambda\mu}.
\end{equation}
Vieta's relations give $ab=\mu/\lambda$ and
$a+b=(\lambda+\mu+\delta)/\lambda$, and the identity that will do most of
the algebraic work in this chapter is
\begin{equation}
AB=(a-1)(1-b)=(a+b)-ab-1=\frac{\delta}{\lambda},
\qquad A+B=a-b.
\label{eq:AB}
\end{equation}
That $b<1<a$ always is immediate from
$\lambda(1-a)(1-b)=\lambda\lrb{1-(a+b)+ab}=-\delta<0$, which places $1$
strictly between the two roots.

\subsection{Three fixation functions}
\label{sec:fixation-funcs}

There are three absorbing-behaviour probabilities, not two, and it pays to
name all of them:
\begin{definition}
\begin{equation}
I(t)=\Prob{\text{no burst by }t},\qquad
D(t)=\Prob{\text{internal extinction by }t,\ \text{no burst}},
\end{equation}
\begin{equation}
\Ifix(t)=\Prob{\xt\notin\{0,H\}}=I(t)-D(t),
\end{equation}
the last being the probability that the cell is still productively infected.
\end{definition}

Both $I$ and $D$ satisfy the same Riccati equation
\begin{equation}
\ddt{I}=\lambda(I-a)(I-b)=\mu-(\lambda+\mu+\delta)I+\lambda I^2,
\label{eq:riccati}
\end{equation}
differently initialised ($I(0)=1$, $D(0)=0$); the branching term $I^2$ comes
from the two independent daughter lineages that a birth of the founder
creates. Solving by separation of variables and writing
$w:=\ee^{\theta t}\ge1$ gives the boxed triplet
\begin{equation}
\boxed{I(t)=\frac{aB+bAw}{B+Aw},\qquad
D(t)=\frac{ab\,(w-1)}{aw-b},\qquad
\Ifix(t)=\frac{(a-b)^2\,w}{(B+Aw)(aw-b)}.}
\label{eq:IDIhat}
\end{equation}
The limits are $I(\infty)=D(\infty)=b$, $\Ifix(\infty)=0$, with
$\Ifix(0)=1$ and $\Ifix'(0)=-(\mu+\delta)$ --- at the instant of founding,
the productive state is lost only by death or catastrophe of the single
founder. An independent confirmation of the third formula arrives in
\S\ref{sec:state-probs}: the state probabilities $p_n(t)$ sum over $n\ge1$
to exactly the right-hand side of \eqref{eq:IDIhat}.

\subsection{Moments and mean release}
\label{sec:moments-recap}

The hazard of bursting at time $t$ is $\delta$ times the current load, so
$I'(t)=-\delta J(t)$; inserting the Riccati equation \eqref{eq:riccati}
then gives the mean
\begin{equation}
\boxed{J(t)=\mean{\xt}=-\frac{1}{\delta}\ddt{I}
=\frac{(a-b)^2\,w}{(B+Aw)^2}
=-\frac{\lambda}{\delta}(I-a)(I-b),}
\label{eq:J}
\end{equation}
with $J(0)=1$ now manifest. The second moment,
\begin{equation}
\boxed{K(t)=\mean{\xt^2}=\Bigl[1+\tfrac{2\lambda}{\delta}(1-I)\Bigr]J
=\frac{2\lambda}{\delta}(\kappa-I)J,\qquad \kappa=1+\frac{\delta}{2\lambda},}
\label{eq:K}
\end{equation}
follows from the moment hierarchy: one checks by differentiation, using
\eqref{eq:riccati}, that the right-hand side satisfies
$K'=2(\lambda-\mu)K+(\lambda+\mu)J-\delta\mean{\xt^3}$ with the third moment
eliminated exactly as in Chapter~3.

The cumulative mean release $V(t)=\mean{\wt}$ increases only at rupture
events, and a rupture from state $n$ occurs at rate $\delta n$ and releases
$n$ particles, so the expected release flux is $\delta\mean{\xt^2}$:
\begin{equation}
V'(t)=\delta K(t),\qquad V(0)=0.
\end{equation}
Integrating with \eqref{eq:J}--\eqref{eq:K} yields the simplification
\begin{equation}
\boxed{V(t)=(1-I)\Bigl[1+\tfrac{\lambda}{\delta}(1-I)\Bigr],}
\label{eq:V}
\end{equation}
a quadratic in $(1-I)$ alone --- no $J$, no separate $\mu$. The equivalent
Chapter~3 form $V=\frac{\lambda-\mu}{\delta}(1-I)-J+1$ reduces to this one
upon substituting $J=-(\lambda/\delta)(I-a)(I-b)$ and using \eqref{eq:AB}.

The lifetime mean release and the conditional mean burst size are
\begin{equation}
\boxed{V_\infty=\frac{a(1-b)}{a-1}=\frac{aB}{A}
=\frac{\lambda-\mu}{\delta}(1-b)+1,
\qquad
\mean{\Kb\mid\text{burst}}=\frac{V_\infty}{1-b}=\frac{a}{a-1}.}
\label{eq:Vinf}
\end{equation}
A small point with large consequences: $V_\infty$ is
$\mean{\Kb\mathbf 1\{\text{burst}\}}$, and so \emph{includes} the
realisations that cleared internally and released nothing (a mass $b$ at
$0$). The conditional mean removes them. When $\mu=0$: $b=0$,
$a=1+\delta/\lambda$, and $V_\infty=1+\lambda/\delta$ coincides with the
conditional mean --- a coincidence that, as we shall see, one should not
lean on.

\subsection{Second moment of the release}
\label{sec:EW2-sec}

For the assay predictions of the discussion we will need $\mean{\wt^2}$.
The release variable jumps once, from $0$ to $X_{\tau^-}$, so
$\ddt{}\mean{\wt^2}=\delta\mean{\xt^3}$. The third moment is eliminated via
the second-moment ODE,
$K'=2(\lambda-\mu)K+(\lambda+\mu)J-\delta\mean{\xt^3}$, giving
\[
\ddt{}\mean{\wt^2}=2(\lambda-\mu)K+(\lambda+\mu)J-K'
=\ddt{}\lrb{\frac{2(\lambda-\mu)}{\delta}V-\frac{\lambda+\mu}{\delta}I-K},
\]
In the second equality we used $V'=\delta K$ together with
$J=-\delta^{-1}I'$. Integrating from $0$ with
$(I,J,K,V,\mean{W^2})(0)=(1,1,1,0,0)$ --- in particular $K(0)=\mean{X_0^2}=1$ ---
\begin{equation}
\boxed{\mean{\wt^2}
=\frac{2(\lambda-\mu)}{\delta}V
-\frac{\lambda+\mu}{\delta}I
-K
+\frac{\lambda+\mu}{\delta}
+1,}
\label{eq:EW2}
\end{equation}
and $\mathrm{Var}(\wt)=\mean{\wt^2}-V^2$. As a check, take $\mu=0$ and
$t\to\infty$: then $I,K\to0$ and $V\to1+\lambda/\delta$, so \eqref{eq:EW2}
gives $\mean{W_\infty^2}=2\lambda^2/\delta^2+3\lambda/\delta+1$, which is
exactly the second moment of the geometric burst-size law of
\S\ref{sec:burst-size} with success probability
$s=\delta/(\lambda+\delta)$, whose second moment is $(2-s)/s^2$.
